\documentclass[aspectratio=169,11pt]{beamer}
\usepackage[T1]{fontenc}
\usepackage{lmodern}
\usepackage[sfdefault,lining,tabular]{FiraSans}
\usepackage[scale=0.85]{FiraMono}
\usepackage{microtype}
\usepackage{array}
\usepackage{booktabs}
\usepackage{tabularx}
\usepackage{colortbl}
\usepackage{listings}
\usepackage{tikz}
\usetikzlibrary{arrows.meta,positioning,calc,shapes.geometric}

\usetheme[sectionpage=none,progressbar=frametitle,numbering=none]{metropolis}

% Use a brighter blue for progress and darker shades for small text.
\definecolor{Navy}{HTML}{10243D}
\definecolor{Ink}{HTML}{172B43}
\definecolor{Muted}{HTML}{45566B}
\definecolor{Blue}{HTML}{174A8B}
\definecolor{Progress}{HTML}{285FA3}
\definecolor{BlueTint}{HTML}{E7F0FC}
\definecolor{Violet}{HTML}{59408B}
\definecolor{VioletTint}{HTML}{F1ECFA}
\definecolor{CodeBg}{HTML}{EFF4FA}
\definecolor{Border}{HTML}{CBD6E4}
\definecolor{RowTint}{HTML}{F0F4F8}

\setbeamercolor{normal text}{fg=Ink,bg=white}
\setbeamercolor{frametitle}{fg=white,bg=Navy}
\setbeamercolor{structure}{fg=Blue}
\setbeamercolor{progress bar}{fg=Progress,bg=Border}
\setbeamercolor{footline}{fg=Muted,bg=white}

\setbeamerfont{frametitle}{size=\large,series=\bfseries}
\setbeamersize{text margin left=1cm,text margin right=1cm}

\setbeamertemplate{itemize item}{\color{Blue}\rule[0.2ex]{0.65ex}{0.65ex}}
\setlength{\leftmargini}{1.1em}

\makeatletter
\setlength{\metropolis@progressinheadfoot@linewidth}{1pt}
\makeatother

\setbeamertemplate{footline}{%
    \begin{beamercolorbox}[wd=\paperwidth,leftskip=1cm,rightskip=1cm]{footline}%
        \color{Border}\rule{\dimexpr\paperwidth-2cm\relax}{0.4pt}\par\vspace{2pt}%
        {\scriptsize\color{Muted}\secname}%
        \hfill%
        {\scriptsize\color{Muted}\insertframenumber\,/\,\inserttotalframenumber}%
        \vspace{4pt}%
    \end{beamercolorbox}%
}

\setbeameroption{hide notes}

% Code listings style
\lstset{
    language=C,
    basicstyle=\fontsize{6.2}{7.8}\selectfont\ttfamily\color{Ink},
    keywordstyle=\color{Blue}\bfseries,
    commentstyle=\color{Muted}\itshape,
    stringstyle=\color{Violet},
    backgroundcolor=\color{CodeBg},
    frame=leftline,
    framerule=2.5pt,
    rulecolor=\color{Ink!40},
    framesep=5pt,
    xleftmargin=7.5pt,
    aboveskip=1pt,
    belowskip=2pt,
    columns=fullflexible,
    keepspaces=true,
    showstringspaces=false,
    escapeinside={(*@}{@*)}
}

% Diagram styles
\tikzset{
    stage/.style={
        draw=Blue,
        fill=BlueTint,
        line width=0.6pt,
        rounded corners=2pt,
        align=center,
        minimum height=8.5mm,
        inner xsep=5pt,
        inner ysep=3pt,
        font=\footnotesize\color{Ink}
    },
    appstage/.style={
        stage,
        draw=Ink!50,
        fill=RowTint
    },
    rejected/.style={
        stage,
        draw=Violet,
        fill=VioletTint
    },
    transition/.style={
        -{Stealth[length=1.8mm]},
        semithick,
        draw=Ink
    }
}

% Reusable helpers
\newcommand{\code}[1]{\texttt{#1}}
\newcommand{\bugline}[1]{\textcolor{Violet}{\bfseries #1}}
\newcommand{\extlink}[2]{\href{#1}{\textcolor{Blue}{#2}}}
\newcommand{\pinnedcode}[2]{\extlink{https://github.com/espressif/esp-idf/blob/7ad199d14a5d0fd2f61193c0f8a51a5d2bfa2c00/components/esp_https_ota/src/esp_https_ota.c#1}{#2}}
\newcommand{\pinnedhdr}[2]{\extlink{https://github.com/espressif/esp-idf/blob/7ad199d14a5d0fd2f61193c0f8a51a5d2bfa2c00/components/esp_https_ota/include/esp_https_ota.h#1}{#2}}
\newcommand{\examplecode}[2]{\extlink{https://github.com/espressif/esp-idf/blob/07bfc09d0c8837b4d149e1151c57635a0c02755a/examples/system/ota/simple_ota_example/main/simple_ota_example.c#1}{#2}}

\newcommand{\passpill}[1]{\colorbox{BlueTint}{\textcolor{Blue}{\bfseries\footnotesize\,#1\,}}}
\newcommand{\failpill}[1]{\colorbox{VioletTint}{\textcolor{Violet}{\bfseries\footnotesize\,#1\,}}}

\newcommand{\srcbar}[2]{%
    \noindent{\footnotesize\ttfamily\color{Ink}#1}\hfill{\scriptsize\color{Muted}#2}\par\vspace{2pt}%
}

\newcommand{\accentcard}[4]{%
    \par\noindent\begin{tikzpicture}[baseline=(card.north)]
    \node[
        fill=#2,
        rounded corners=2pt,
        inner xsep=9pt,
        inner ysep=5pt,
        text width=\dimexpr\linewidth-18pt\relax,
        align=flush left,
        path picture={
            \fill[#1] (path picture bounding box.south west) rectangle
            ([xshift=3pt]path picture bounding box.north west);
        }
    ] (card) {%
            {\small\bfseries\color{#1}#3}\par\vspace{2pt}%
        {\footnotesize\color{Ink}\setlength{\parskip}{2pt}#4}%
    };
    \end{tikzpicture}\par%
}

\newcommand{\pillarcard}[2]{%
    \noindent\begin{tikzpicture}[baseline=(card.north)]
    \node[
        fill=RowTint,
        draw=Border,
        line width=0.5pt,
        rounded corners=2pt,
        inner xsep=9pt,
        inner ysep=8pt,
        minimum height=3.25cm,
        text width=\dimexpr\linewidth-18.5pt\relax,
        align=flush left,
        path picture={
            \fill[Blue] (path picture bounding box.north west) rectangle
            ([yshift=-3pt]path picture bounding box.north east);
        }
    ] (card) {%
        \parbox[t][2.65cm][t]{\linewidth}{%
            \rightskip=0pt plus 1fil\hyphenpenalty=10000\exhyphenpenalty=10000
            {\normalsize\bfseries\color{Blue}#1}\par\vspace{5pt}%
            {\footnotesize\color{Ink}#2}%
        }%
    };
    \end{tikzpicture}%
}

\begin{document}

% Slide 1: Table of contents
    \begin{frame}{The Technical Case (Table of Contents)}
        \begingroup
        \renewcommand{\arraystretch}{1.45}
        \small
        \begin{tabularx}{\textwidth}{@{}l >{\raggedright\arraybackslash}X r@{}}
            \toprule
            & \textbf{Questions (from the main README)} & \textbf{Slides} \\
            \midrule
            Q1 & Are there bugs in the OTA implementation? & \hyperlink{ota-bugs}{\textcolor{Blue}{3--8}} \\
            Q2 & What can go wrong when calling the \code{esp\_https\_ota} wrapper as it's implemented now? & \hyperlink{ota-bugs}{\textcolor{Blue}{3--6}}, \hyperlink{ota-example}{\textcolor{Blue}{9}} \\
            Q3 & How should we use \code{esp\_https\_ota} to ensure sufficient robustness? Is it even possible? & \hyperlink{ota-example}{\textcolor{Blue}{9--11}} \\
            Q4 & How would we make sure that OTA in the product is robust enough? & \hyperlink{ota-hardware}{\textcolor{Blue}{12--15}} \\
            \bottomrule
        \end{tabularx}
        \endgroup
        \vspace{3mm}

        {\footnotesize\color{Muted} The OTA overview is on slide 2, supplementary materials are on slide 16, and LLM use is on slide 17.}
    \end{frame}


    \section{What can go wrong?}

% Slide 2: Overview flow
    \begin{frame}{The OTA update stages (Preamble)}
        \begin{center}
            \begin{tikzpicture}[node distance=5mm]
                \node[stage] (download) {Download\\and write};
                \node[stage,right=of download] (validate) {Check\\image};
                \node[stage,right=of validate] (select) {Select\\boot slot};
                \node[appstage,right=of select] (reboot) {Reboot};
                \node[appstage,right=of reboot] (confirm) {Confirm\\health};
                \draw[transition] (download) -- (validate);
                \draw[transition] (validate) -- (select);
                \draw[transition] (select) -- (reboot);
                \draw[transition] (reboot) -- (confirm);
                \draw[Blue,semithick] (download.south west) -- ++(0,-2mm) -| (select.south east);
                \node[below=2.5mm,font=\footnotesize] at ($(download.south west)!0.5!(select.south east)$)
                    {\pinnedcode{\#L381}{\code{esp\_https\_ota()}}};
                \draw[Ink!60,semithick] (reboot.south west) -- ++(0,-2mm) -| (confirm.south east);
                \node[below=2.5mm,font=\footnotesize,color=Ink] at ($(reboot.south west)!0.5!(confirm.south east)$)
                    {Charger firmware};
            \end{tikzpicture}
        \end{center}
        \vspace{3mm}
        \noindent
        \begin{minipage}[t]{0.485\textwidth}
            \accentcard{Blue}{BlueTint}{What the OTA component does}{%
                The wrapper function (\pinnedcode{\#L381}{\code{esp\_https\_ota()}}) downloads the image, checks it, and selects the next boot slot.%
            }
        \end{minipage}\hfill
        \begin{minipage}[t]{0.485\textwidth}
            \accentcard{Ink}{RowTint}{What the charger does}{%
                In the example \examplecode{\#L106-L114}{\code{simple\_ota\_example.c}}, a reboot happens as soon as \code{ESP\_OK} is returned. The charger then must decide if the reboot is safe and run first-boot checks.%
            }
        \end{minipage}
        \note{Target time is 60 seconds. Show what the API covers and what the product owns. The example restarts right away when the wrapper returns success. Safe charging pauses and saved data safety belong to the charger firmware.}
    \end{frame}

% Slide 3: Bug 1
    \begin{frame}[fragile,label=ota-bugs]{Boot selection can fail, but \code{ESP\_OK} is returned (Bug\#1)}
        \begin{columns}[T,onlytextwidth]
            \begin{column}{0.53\textwidth}
                \srcbar{\code{esp\_https\_ota\_finish()}}{\pinnedcode{\#L340-L366}{lines 340 to 366}}
                \begin{lstlisting}
esp_err_t err = ESP_OK;
// Line 344 sets outer err inside switch.
if ((err == ESP_OK) &&
    (handle->state == ESP_HTTPS_OTA_SUCCESS)) {
    (*@\bugline{esp\_err\_t err = esp\_ota\_set\_boot\_partition(}@*)
        (*@\bugline{handle->update\_partition);}@*)
    // Line 360 hides outer err; outer stays 0.
}
return err;
                \end{lstlisting}
            \end{column}
            \begin{column}{0.44\textwidth}
                \accentcard{Violet}{VioletTint}{Injected boot failure}{%
                    Boot selection returned \code{ESP\_ERR\_FLASH\_OP\_FAIL} (\code{24577}). The OTA call returned \code{ESP\_OK} (\code{0}).%
                }
                \vspace{2mm}
                \accentcard{Ink}{RowTint}{Why this happens}{%
                    An inner \code{err} variable shadows the outer error variable at \pinnedcode{\#L360}{line 360}. The function logs the error and returns \code{ESP\_OK}.\par
                    Issue \extlink{https://github.com/espressif/esp-idf/issues/12635}{\#12635} reports this bug. The \extlink{https://github.com/espressif/esp-idf/commit/d778b07b2213370e1603d1aadf7b60393b3f7569}{v5.1.3} fix updates the outer variable.%
                }
            \end{column}
        \end{columns}
        \note{Target time is 2 minutes. Point out the outer and inner err variables in esp\_https\_ota\_finish. The fix writes to the outer variable. Notice esp\_ota\_set\_boot\_partition validates the image a second time, so a bootloader\_mmap failure under memory or MMU pressure (issue 16968) can trigger this bug in practice. The injected test shows false success. Physical flash failure and next-boot behavior on real hardware are still untested. Upstream issue 12635 reports the same bug.}
    \end{frame}

% Slide 4: Bug 2
    \begin{frame}[fragile]{Setup loses the exact error (Bug\#2)}
        \begin{columns}[T,onlytextwidth]
            \begin{column}{0.55\textwidth}
                \srcbar{\code{esp\_https\_ota()}}{\pinnedcode{\#L392-L396}{lines 392 to 396}}
                \begin{lstlisting}
esp_https_ota_handle_t https_ota_handle = NULL;
esp_err_t err = esp_https_ota_begin(
    &ota_config, &https_ota_handle);
if (https_ota_handle == NULL) {
    (*@\bugline{return ESP\_FAIL;}@*) // Drops err at line 395
}
                \end{lstlisting}
            \end{column}
            \begin{column}{0.42\textwidth}
                \accentcard{Violet}{VioletTint}{Injected memory failure}{%
                    The test expected \code{ESP\_ERR\_NO\_MEM} (\code{257}). The wrapper function returned \code{ESP\_FAIL} (\code{-1}).%
                }
                \vspace{2.5mm}
                \accentcard{Ink}{RowTint}{Why this matters}{%
                    The header \pinnedhdr{\#L53-L60}{\code{esp\_https\_ota.h}} (lines 53 to 60) documents specific error codes.\par
                    The wrapper gives the retry logic a generic failure code when setup fails.%
                }
            \end{column}
        \end{columns}
        \note{Target time is 75 seconds. Look at esp\_https\_ota right after begin fails. Returning the exact error helps the caller pick a smart retry rule. The test forces calloc inside the component to fail.}
    \end{frame}

% Slide 5: Bug 3
    \begin{frame}[fragile]{HTTP errors can start a flash update (Bug\#3)}
        \begin{center}
            \begin{tikzpicture}[node distance=5mm]
                \node[rejected] (http) {HTTP 404\\or 500};
                \node[appstage,right=of http] (accepted) {Connection\\accepted};
                \node[stage,right=of accepted] (flash) {\code{esp\_ota\_begin()}\\erases slot};
                \node[rejected,right=of flash] (validate) {Image check fails\\(\code{5379})};
                \draw[transition] (http) -- (accepted);
                \draw[transition] (accepted) -- (flash);
                \draw[transition] (flash) -- (validate);
            \end{tikzpicture}
        \end{center}
        \vspace{3mm}
        \begin{columns}[T,onlytextwidth]
            \begin{column}{0.56\textwidth}
                \srcbar{\code{\_http\_handle\_response\_code()}}{\pinnedcode{\#L60-L88}{lines 60 to 88, 271}}
                \begin{lstlisting}
// Redirect and 401 handling omitted here.
if (process_again(status_code)) {
    // Comment says redirects only (line 74),
    // yet process_again includes 401 (line 52).
}
(*@\bugline{return ESP\_OK;}@*) // Reached for HTTP 404 and 500
                \end{lstlisting}
            \end{column}
            \begin{column}{0.41\textwidth}
                \accentcard{Violet}{VioletTint}{Flash wear and endless retries}{%
                    HTTP 404 and 500 must fail before any flash erase, with zero calls to \code{esp\_ota\_begin()}.\par
                    HTTP 401 also retries with no limit inside \pinnedcode{\#L109}{\code{\_http\_connect()}}.%
                }
            \end{column}
        \end{columns}
        \note{Target time is 105 seconds. Show \_http\_handle\_response\_code and the first branch of esp\_https\_ota\_perform. Note that esp\_ota\_begin with OTA\_SIZE\_UNKNOWN erases the passive slot before checking the HTTP body. Also note that \_http\_connect retries HTTP 401 with no limit and does not enforce redirect limits in either 7ad199d14a or v6.1.1. Checking only the final error would hide the flash erase.}
    \end{frame}

% Slide 6: Bug 4
    \begin{frame}[fragile]{The wrapper has no overall time limit (Bug\#4)}
        \begin{columns}[T,onlytextwidth]
            \begin{column}{0.54\textwidth}
                \srcbar{\code{esp\_https\_ota\_perform()}}{\pinnedcode{\#L285-L306}{lines 285 to 306, 398 to 403}}
                \begin{lstlisting}[basicstyle=\fontsize{6.1}{7.6}\selectfont\ttfamily\color{Ink}]
data_read = esp_http_client_read(
    handle->http_client, ...);
if (data_read == 0) {
    bool done =
        esp_https_ota_is_complete_data_received(
            https_ota_handle);
    if ((errno == ENOTCONN || errno == ECONNRESET ||
         errno == ECONNABORTED) && !done) {
        return ESP_FAIL;
    }
    if (!done) {
        (*@\bugline{return ESP\_ERR\_HTTPS\_OTA\_IN\_PROGRESS;}@*)
    }
}
                \end{lstlisting}
            \end{column}
            \begin{column}{0.43\textwidth}
                \accentcard{Violet}{VioletTint}{Loop with no limit}{%
                    A read timeout leaves \code{errno} at zero in \code{7ad199d14a} and returns \code{-ESP\_ERR\_HTTP\_EAGAIN} in \code{v6.1.1}.\par
                    \pinnedcode{\#L398-L403}{\code{esp\_https\_ota()}} loops in \code{while (1)} with no deadline.%
                }
                \vspace{2.5mm}
                \accentcard{Ink}{RowTint}{QEMU stall test}{%
                    Both unpatched builds retry until the server closes the connection. The local no-progress guard returns \code{ESP\_ERR\_TIMEOUT} (\code{263}).%
                }
            \end{column}
        \end{columns}
        \note{Target time is 75 seconds. Walk through esp\_https\_ota\_perform at lines 285 to 306 and the while loop in esp\_https\_ota at lines 398 to 403. In QEMU, a stalled 200 stream triggers 14 read timeouts on 7ad199d14a and 15 on v6.1.1 until the test server closes after 8 seconds. The local fix stops with ESP\_ERR\_TIMEOUT (263) after 4 timeouts.}
    \end{frame}

% Slide 7: QEMU results
    \begin{frame}{One baseline and five tests (to reproduce the bugs)}
        \begingroup
        \setlength{\aboverulesep}{0pt}
        \setlength{\belowrulesep}{0pt}
        \renewcommand{\arraystretch}{1.05}
        \footnotesize
        \begin{tabularx}{\textwidth}{ X c c c }
    \toprule
    \rowcolor{Ink!10}
    \textbf{Case}                      & \textbf{Assigned SDK}                             & \textbf{Local fixes} & \textbf{v6.1.1}             \\
    \midrule
    Good download and image check      & \passpill{Pass} (\code{0})                        & \passpill{Pass}      & \passpill{Pass}             \\
    \rowcolor{RowTint}
    Boot selection fails (injected)    & \failpill{Fail} (\code{0}, expected \code{24577}) & \passpill{Pass} & \passpill{Pass} \\
    Memory allocation fails (injected) & \failpill{Fail} (\code{-1}, expected \code{257}) & \passpill{Pass} & \passpill{Pass} \\
    \rowcolor{RowTint}
    HTTP 404 response                  & \failpill{Fail} (\code{5379}, \code{begin=1})     & \passpill{Pass} & \passpill{Pass} \\
    HTTP 500 response                  & \failpill{Fail} (\code{5379}, \code{begin=1})     & \passpill{Pass} & \passpill{Pass} \\
    \rowcolor{RowTint}
    Stalled download (no deadline)     & \failpill{Fail} (\code{-1}, expected \code{263})  & \passpill{Pass} & \failpill{Fail} (\code{-1}) \\
    \bottomrule
        \end{tabularx}
        \endgroup
        \par\vspace{3mm}
        \accentcard{Ink}{RowTint}{Test limits}{%
            Single-core ESP32 QEMU over HTTP with Timer Group watchdogs disabled. \code{begin=1} means one call to \code{esp\_ota\_begin()}.%
        }
        \note{Target time is 90 seconds. A language model flagged paths, which I checked against the pinned SDK and QEMU results. The baseline stops after boot-slot selection. The two unpatched builds retry the stalled download until the server closes the connection.}
    \end{frame}

% Slide 8: Bug summary table
    \begin{frame}{Bugs, fixes, and upstream reports}
        \begingroup
        \setlength{\aboverulesep}{0pt}
        \setlength{\belowrulesep}{0pt}
        \renewcommand{\arraystretch}{1.12}
        \footnotesize
        \begin{tabularx}{\textwidth}{
                >{\raggedright\arraybackslash}p{0.31\textwidth}
                >{\raggedright\arraybackslash}p{0.22\textwidth}
                >{\raggedright\arraybackslash}p{0.18\textwidth}
                >{\raggedright\arraybackslash}X }
    \toprule
    \rowcolor{Ink!10}
    \textbf{Bug} & \textbf{Pinned code} & \textbf{Fix source} & \textbf{ESP-IDF issue} \\
    \midrule
    False success on boot selection
    & \pinnedcode{\#L360}{Line 360}
    & \extlink{https://github.com/espressif/esp-idf/commit/d778b07b2213370e1603d1aadf7b60393b3f7569}{v5.1.3}
    & \extlink{https://github.com/espressif/esp-idf/issues/12635}{\#12635} (exact) \\
    \rowcolor{RowTint}
    Setup error lost
    & \pinnedcode{\#L393-L395}{Lines 393 to 395}
    & \extlink{https://github.com/espressif/esp-idf/blob/v6.1.1/components/esp_https_ota/src/esp_https_ota.c\#L1099-L1102}{v6.1.1}
    & No exact report \\
    Flash erase on HTTP 404 or 500
    & \pinnedcode{\#L60-L110}{Lines 60 to 110}, \pinnedcode{\#L271}{271}
    & \extlink{https://github.com/espressif/esp-idf/blob/v4.4.8/components/esp_https_ota/src/esp_https_ota.c}{v4.4.8}
    & \extlink{https://github.com/espressif/esp-idf/issues/7058}{\#7058}, \extlink{https://github.com/espressif/esp-idf/issues/8195}{\#8195} (related) \\
    \rowcolor{RowTint}
    No overall time limit for OTA
    & \pinnedcode{\#L285-L306}{Lines 285 to 306}, \pinnedcode{\#L398-L403}{398 to 403}
    & Not fixed yet (in \code{v6.1.1})
    & \extlink{https://github.com/espressif/esp-idf/issues/6364}{\#6364} (related) \\
    \bottomrule
        \end{tabularx}
        \endgroup
        \note{Target time is 45 seconds. The fix source column points to code. Slide 7 shows which builds ran in QEMU.}
    \end{frame}


    \section{How should the application call it?}

% Slide 9: Example analysis
    \begin{frame}[fragile,label=ota-example]{Why the code in \code{simple\_ota\_example.c} is not reliable enough?}
        \begin{columns}[T,onlytextwidth]
            \begin{column}{0.52\textwidth}
                \srcbar{\code{simple\_ota\_example.c}}{\examplecode{\#L106-L114}{lines 106 to 114}}
                \begin{lstlisting}[basicstyle=\fontsize{6.2}{7.8}\selectfont\ttfamily\color{Ink}]
esp_err_t ret = esp_https_ota(&config);
if (ret == ESP_OK) {
    (*@\bugline{esp\_restart();}@*) // No charging-state check
} else {
    ESP_LOGE(TAG, "Firmware upgrade failed");
}
(*@\bugline{while (1) \{}@*) // Stuck task; no retry
    (*@\bugline{vTaskDelay(1000 / portTICK\_PERIOD\_MS);}@*)
(*@\bugline{\}}@*)
                \end{lstlisting}
            \end{column}
            \begin{column}{0.45\textwidth}
                \accentcard{Ink}{RowTint}{Problems in the example code}{%
                    The example never confirms the new image. Startup can also \examplecode{\#L144-L152}{erase NVS} (including the credentials stored there).%
                }
                \vspace{2.5mm}
                \accentcard{Violet}{VioletTint}{Can the wrapper be safe enough?}{%
                    No. The one-shot wrapper hides boot and setup errors, skips checks before flash erase, and has no overall time limit.%
                }
            \end{column}
        \end{columns}
        \note{Target time is 75 seconds. Answer questions 2 and 3 from the brief. The example restarts on reported success without checking whether the charger is active.}
    \end{frame}

% Slide 10: Staged API sketch
    \begin{frame}[fragile]{Step-by-step calls give more control}
        \begin{columns}[T,onlytextwidth]
            \begin{column}{0.55\textwidth}
                \srcbar{Step-by-step caller sketch}{\pinnedcode{\#L210}{lines 210, 293}}
                \begin{lstlisting}[basicstyle=\fontsize{6.1}{7.6}\selectfont\ttfamily\color{Ink}]
err = esp_https_ota_begin(&cfg, &handle);
if (err != ESP_OK) return err;
err = esp_https_ota_get_img_desc(handle, &info);
if (err == ESP_OK && !valid_header(&info))
    err = ESP_ERR_INVALID_VERSION;
if (err == ESP_OK) do {
    err = timed_out() ? ESP_ERR_TIMEOUT
        : esp_https_ota_perform(handle);
} while (err == ESP_ERR_HTTPS_OTA_IN_PROGRESS);
// Use abort() on error if the SDK provides it.
esp_err_t finish_err = esp_https_ota_finish(handle);
return err == ESP_OK ? finish_err : err;
                \end{lstlisting}
            \end{column}
            \begin{column}{0.42\textwidth}
                \accentcard{Blue}{BlueTint}{Check before erase}{%
                    Call \code{esp\_https\_ota\_get\_img\_desc()} before any flash erase.%
                }
                \vspace{2mm}
                \accentcard{Ink}{RowTint}{Time-limited loop}{%
                    Run in one worker task and check the deadline between calls.%
                }
                \vspace{2mm}
                \accentcard{Ink}{RowTint}{Safe reboot}{%
                    Reboot only after finish succeeds and charging is safe.%
                }
            \end{column}
        \end{columns}
        \note{Target time is 105 seconds. The sketch uses finish for cleanup on 7ad199d14a. Use abort on error when the SDK provides it. The application must also serialize OTA (issue 18127) and keep the OTA task stack in internal RAM (issue 10437).}
    \end{frame}

% Slide 11: Staged API limitations
    \begin{frame}{The caller still needs a fixed SDK}
        \begin{columns}[T,onlytextwidth]
            \begin{column}{0.50\textwidth}
                \accentcard{Ink}{RowTint}{Limits inside the pinned SDK}{%
                    The step-by-step API uses the same HTTP and boot-selection code as \code{esp\_https\_ota()}.\par
                    Code at \code{7ad199d14a} also allows HTTPS downgrades (\extlink{https://github.com/espressif/esp-idf/commit/eef3b0bf}{eef3b0bf}) and blocks certificate bundles (\extlink{https://github.com/espressif/esp-idf/issues/5504}{\#5504}).\par
                    \code{begin()} and \pinnedcode{\#L231-L245}{\code{get\_img\_desc()}} can block in read loops, where caller timers cannot stop them.%
                }
            \end{column}
            \begin{column}{0.47\textwidth}
                \accentcard{Blue}{BlueTint}{What must change}{%
                    Use a fixed SDK or backport the fixes. Limit network waits and retries. Then test the overall time limit.\par
                    Even newer SDKs need caller guards for redirect limits and open \code{esp\_ota\_ops} bugs.%
                }
                \vspace{2.5mm}
                \accentcard{Ink}{RowTint}{Retry rules}{%
                    Retry short network errors with backoff. Stop right away on bad certificates or bad images.%
                }
            \end{column}
        \end{columns}
        \note{Target time is 90 seconds. One worker task keeps the charger responsive during downloads. Beyond the four core bugs, 7ad199d14a allows HTTPS downgrades (commit eef3b0bf) and blocks certificate bundles (issue 5504). Newer SDKs such as v6.1.1 still have unguarded redirect loops in \_http\_connect and open esp\_ota\_ops bugs (issues 19413 and 19439).}
    \end{frame}


    \section{How do we test charger OTA?}

% Slide 12: Hardware verification matrix
    \begin{frame}[label=ota-hardware]{Test recovery on real hardware}
        \begingroup
        \setlength{\aboverulesep}{0pt}
        \setlength{\belowrulesep}{0pt}
        \renewcommand{\arraystretch}{1.12}
        \footnotesize
        \begin{tabularx}{\textwidth}{ >{\raggedright\arraybackslash}p{0.39\textwidth} >{\raggedright\arraybackslash}X }
    \toprule
    \rowcolor{Ink!10}
    \textbf{Fault case} & \textbf{Expected result}                                                   \\
    \midrule
    Network stalls, drops, and cut-off downloads
    & The update stops in the time limit. The old firmware still works.      \\
    \rowcolor{RowTint}
    Bad certificate or unapproved image
    & The update is rejected. Current firmware stays usable.                     \\
    Power loss during flash write, boot selection, or first boot
    & A working image boots, the charger stays safe, and saved settings survive. \\
    \rowcolor{RowTint}
    Failed first-boot health check
    & The charger rolls back, and old firmware keeps its data.                   \\
    \bottomrule
        \end{tabularx}
        \endgroup
        \vspace{2.5mm}
        \accentcard{Blue}{BlueTint}{Before testing}{%
            Use production TLS, watchdog, and bootloader settings; test real heap load (\extlink{https://github.com/espressif/esp-idf/issues/16968}{\#16968}) and older field builds (\extlink{https://github.com/espressif/esp-idf/issues/9987}{\#9987}, \extlink{https://github.com/espressif/esp-idf/issues/17595}{\#17595}).%
        }
        \note{Target time is 105 seconds. Describe a hardware bench with network faults, power cuts, and charging checks. Include dual-core flash watchdog checks (issues 10890 and 14450), realistic firmware sizes under heap and MMU load (issue 16968), and upgrades from older field firmware (issues 9987 and 17595). Agree on update time limits, charging pause rules, and recovery goals before testing. QEMU disables watchdogs and runs one core over HTTP.}
    \end{frame}

% Slide 13: Rollback flow
    \begin{frame}{A failed first boot must roll back}
        \begin{center}
            \begin{tikzpicture}[node distance=9mm]
                \node[stage] (boot) {Boot new image};
                \node[stage,right=of boot] (health) {Run health tests};
                \coordinate[right=8mm of health.east] (branch);
                \node[stage,anchor=west] (confirm) at ([xshift=29mm,yshift=8.5mm]health.east)
                    {Confirm\\new firmware};
                \node[rejected,anchor=west] (rollback) at ([xshift=29mm,yshift=-8.5mm]health.east)
                    {Roll back to\\last good image};
                \draw[transition] (boot) -- (health);
                \draw[semithick,draw=Ink] (health.east) -- (branch);
                \draw[transition] (branch) |-
                node[pos=0.75,above,font=\footnotesize,color=Blue] {Pass} (confirm.west);
                \draw[transition] (branch) |-
                node[pos=0.75,below,font=\footnotesize,color=Violet] {Fail or reset} (rollback.west);
            \end{tikzpicture}
        \end{center}
        \vspace{1.5mm}
        \noindent
        \begin{minipage}[t]{0.485\textwidth}
            \accentcard{Blue}{BlueTint}{Confirm the new image}{%
                Turn on bootloader rollback. Call\par\vspace{1pt}
                {\fontsize{6.4}{8}\selectfont\code{esp\_ota\_mark\_app\_valid\_cancel\_rollback()}}\par\vspace{1pt}
                only after the checks pass.%
            }
        \end{minipage}\hfill
        \begin{minipage}[t]{0.485\textwidth}
            \accentcard{Violet}{VioletTint}{Roll back to a bootable image}{%
                If checks fail, call\par
                {\fontsize{6.4}{8}\selectfont\code{esp\_ota\_mark\_app\_invalid\_rollback\_and\_reboot()}}.\par
                A reset before confirmation rolls back if a previous image can boot.%
            }
        \end{minipage}
        \note{Target time is 75 seconds. A healthy image calls esp\_ota\_mark\_app\_valid\_cancel\_rollback. A failed test requests rollback. Watch for open issue 19413, where a retry after a failed otadata write can return ESP\_OK without confirming the running firmware. A reset before confirmation also rolls back when a good old image exists. Test this on real hardware.}
    \end{frame}

% Slide 14: Rollout policy
    \begin{frame}{Use a canary/gradual deployment before a full rollout}
        \begin{center}
            \begin{tikzpicture}[node distance=18mm]
                \node[stage] (s1) {\textbf{Stage 1}\\Lab tests};
                \node[stage,right=of s1] (s2) {\textbf{Stage 2}\\Small-group rollout};
                \node[appstage,right=of s2] (s3) {\textbf{Stage 3}\\Full rollout};
                \draw[transition] (s1) -- node[above,font=\scriptsize,color=Blue] {Pass} (s2);
                \draw[transition] (s2) -- node[above,font=\scriptsize,color=Ink] {Healthy} (s3);
            \end{tikzpicture}
        \end{center}
        \vspace{3mm}
        \noindent
        \begin{minipage}[t]{0.485\textwidth}
            \accentcard{Blue}{BlueTint}{What to keep an eye on}{%
                \begin{itemize}
                            \setlength{\itemsep}{1pt}\setlength{\topsep}{2pt}\setlength{\parsep}{0pt}
                        \item Reboot tests during idle and active charging
                        \item Firmware version reported after reboot
                        \item Time to confirmed boot and rollback rate
                        \item Chargers that stop reporting
                \end{itemize}%
            }
        \end{minipage}\hfill
        \begin{minipage}[t]{0.485\textwidth}
            \accentcard{Ink}{RowTint}{When to pause}{%
                Agree on rollback, stall, and offline limits with product and operations teams before a wider rollout.\par
                Pause the rollout right away if the small group crosses any limit.%
            }
        \end{minipage}
        \note{Target time is 45 seconds. Agree on pause limits with product and operations teams. Always read the firmware version from the charger after it reboots.}
    \end{frame}

% Slide 15: How I would ship OTA
    \begin{frame}{How I would ship OTA}
        \vspace{1mm}
        \noindent
        \begin{minipage}[t]{0.315\textwidth}
            \pillarcard{Fixed SDK}{%
                Pick an ESP-IDF version or patched component that passes all the regression tests.%
            }
        \end{minipage}\hfill
        \begin{minipage}[t]{0.315\textwidth}
            \pillarcard{Bounded caller}{%
                Run step-by-step calls in one task with time limits, few retries, and a safe reboot check.%
            }
        \end{minipage}\hfill
        \begin{minipage}[t]{0.315\textwidth}
            \pillarcard{Hardware proof}{%
                Test safe charging, first-boot checks, and remote recovery on real hardware.%
            }
        \end{minipage}
        \note{Target time is 90 seconds. Tie each point back to the four questions in the brief. Note that upgrading ESP-IDF fixes older bugs, yet newer releases still add regressions such as issues 13065 and 19439, so automated fault tests remain necessary. Invite questions on time limits, rollback checks, and fault tests.}
    \end{frame}


    \section{Supplementary materials}

% Slide 16: Source and reproduction
    \begin{frame}{Supplementary materials (source code and test instructions)}
        \begin{columns}[T,onlytextwidth]
            \begin{column}{0.50\textwidth}
                \accentcard{Ink}{RowTint}{Code and upstream links}{%
                    \begin{itemize}
                                \setlength{\itemsep}{2pt}\setlength{\topsep}{2pt}\setlength{\parsep}{0pt}
                            \item \pinnedcode{}{OTA source (assigned for review) at \code{7ad199d14a}}
                            \item \examplecode{}{Caller example (assigned for review) at \code{07bfc09d0c}}
                            \item \extlink{https://github.com/espressif/esp-idf/issues/12635}{Boot-selection issue \#12635}
                            \item \extlink{https://github.com/espressif/esp-idf/issues/7058}{HTTP status issues \#7058} and \extlink{https://github.com/espressif/esp-idf/issues/8195}{\#8195}
                            \item \extlink{https://github.com/espressif/esp-idf/issues/6364}{Time limit issue \#6364}
                    \end{itemize}%
                }
            \end{column}
            \begin{column}{0.47\textwidth}
                \accentcard{Blue}{BlueTint}{Run the QEMU tests}{%
                    \colorbox{white}{\parbox{\dimexpr\linewidth-6pt\relax}{\ttfamily\footnotesize\color{Ink}nix develop -c make test test-fixed}}\par\vspace{5pt}
                    Results are in \code{docs/qemu-results.md}.\par
                    More issues are in \code{docs/related-issues.md}.%
                }
            \end{column}
        \end{columns}
        \note{Use this slide during questions and answers for commit hashes, issue numbers, or test commands.}
    \end{frame}

% Slide 17: LLM use
    \begin{frame}{AI/LLM use}
        \begin{itemize}
            \item I use AI-assisted coding tools for finding possible bugs and implementing the various things like the tests and boilerplate.
            \item I use an AI skill called ``Ponytail'' to make the changes and work AI does bounded and small, so I can see everything and make sure AI doesn't make a serious mistake.
            \item I use AI to discuss different aspects of the code (given the code as context). Also, I use AI to search for relevant issues and posts that could help me with the task.
            \item I try to not use AI for writing text because AI writes very badly.
            \item AI is very competent, but it can't take responsibility for what it does. Additionally, I find that current AI models can overengineer and do more than I asked them to do.
            \item I also use (opinionated) templates for my project, which I find help AI work closer to how I myself would've done things.
        \end{itemize}
    \end{frame}

\end{document}
